% Section 4: Temporal Stability — Summary of C.3 and C.4 (~1,100 words)

\paragraph{Stability concepts defined.}
We distinguish two stability concepts throughout this section.
(1)~\textbf{Metric stability}: the mean PP score varies less than $0.010$
across census years for a given state --- a property of the score
distribution. (2)~\textbf{Boundary stability}: the Jaccard similarity
(Intersection over Union, IoU) between district assignments across census
years is $\geq 0.70$ --- a property of the geographic shapes themselves.
C.3 measures metric stability; C.4 measures boundary stability. Both are
necessary for temporal validation: a plan that produces similar PP scores
but places radically different tracts in each district provides no
continuity benefit for voters or incumbents, while a plan whose boundaries
change little but whose compactness drifts substantially may not satisfy
DIA's longitudinal certification standard.

\subsection{Two Distinct Notions of Temporal Stability}

The word ``stability'' covers two different empirical questions in
redistricting research, both important but analytically distinct.

\textbf{Decadal boundary stability (C.3)} asks: when the algorithm is
applied to two successive census years for the same state, how similar
are the resulting district boundaries? Stability in this sense matters
for voters (who may have to learn new districts), incumbents (whose
electoral bases change), and courts (which have expressed preference
for continuity in redistricting challenges). It is a property of the
algorithm's behavior across \emph{time}, holding the state fixed.

\textbf{Historical performance convergence (C.4)} asks: is the gap
between algorithmic plans and enacted plans shrinking over time, as
redistricting reform spreads and enacted maps improve? This is a
property of the algorithm's \emph{relative} performance compared to
a changing benchmark.

Papers C.3 \citep{dellalibera2026c3} and C.4 \citep{dellalibera2026c4}
address these questions respectively.

\subsection{C.3: Geographic Clustering Stability Across Censuses}

C.3 conducts a systematic comparison of recursive bisection versus
n-way partitioning methods for maintaining district continuity across
the 2010--2020 decade. Five southern states with significant minority
populations were selected as the test set: Alabama, Georgia, Louisiana,
Mississippi, and North Carolina. Both algorithms were run on 2010 and
2020 census data (20 total redistricting runs), and decadal stability
was measured using Intersection over Union (IoU) between matched
districts across the two years.

The key finding is that \textbf{recursive bisection provides measurably
better temporal stability}: mean population disruption is $71.6\%$
for recursive bisection versus $72.4\%$ for n-way partitioning, a
statistically significant $0.8$ percentage-point advantage ($p < 0.05$,
paired $t$-test). The advantage appears in 4 of 5 test states.

For a state with 5 million residents, this $0.8$ pp advantage translates
to approximately 40,000 fewer people reassigned to different districts
across the decade---a modest but consistent benefit. The mechanism is
intuitive: hierarchical bisection preserves coarse-level geographic
divisions (north vs.\ south half of a state) that direct $k$-way
optimization may not respect, because the coarse structure is embedded
in the bisection tree.

\subsection{C.3: Within-Slice Temporal Correlation}

The C.2 slice framework applied to C.3's data yields median within-slice
cross-decade IoU of $0.71$ (IQR: $[0.63, 0.79]$). This means that
districts drawn for 2010 and 2020 share on average 71\% of their
geographic area within any given geographic slice. The overlap is higher
in rural slices ($0.78$) and lower in high-growth urban slices ($0.61$),
consistent with the intuition that rapidly growing metropolitan areas
force larger boundary revisions.

District stability is not a goal the algorithm explicitly optimizes---
METIS maximizes edge-cut quality for the current census year without
memory of previous years. The 71\% IoU is therefore a \emph{structural}
consequence of geographic stability: places that are geographically
coherent in 2010 tend to remain so in 2020, guiding the algorithm to
similar partitions.

\subsection{C.4: Longitudinal Convergence of Algorithmic and Enacted Plans}

Paper C.4 \citep{dellalibera2026c4} takes a twenty-year view, examining
compactness trajectories for both algorithmic and enacted plans from
the 2000 through the 2020 redistricting cycles. Three census cycles
provide three data points per state---enough to estimate trend
directions, though not to fit complex models.

\subsubsection{Algorithmic plan compactness is stable}

For the algorithmic plans, national mean Polsby-Popper is $0.412$
(2000), $0.428$ (2010), and $0.441$ (2020). The improvement of
$\Delta PP = +0.029$ over twenty years reflects refinements in input
data quality, not algorithmic changes. Coefficient of variation across
years is $3.5\%$---small relative to cross-state variation ($14.2\%$).
Algorithm behavior is \textbf{temporally stable}.

\subsubsection{Enacted plan compactness is improving, narrowing the gap}

For enacted plans, the mean Polsby-Popper improved from $0.298$ (2000)
to $0.316$ (2010) to $0.331$ (2020)---a $\Delta PP = +0.033$ gain
over twenty years. This improvement reflects the spread of redistricting
reform commissions, court interventions striking down egregiously
non-compact plans (Pennsylvania 2018, North Carolina 2019), and
improvements in mapping software that make moderately compact plans
easier to draw.

The gap between algorithmic and enacted plan compactness therefore
\textbf{narrowed} from $0.412 - 0.298 = 0.114$ (2000) to
$0.441 - 0.331 = 0.110$ (2020). The algorithmic advantage persists
but decreases by $0.004$ per cycle. Extrapolating (with the caveat
that two trend points do not make a reliable forecast), the gap would
close around 2060 if reform continues at the current pace---well beyond
the planning horizon for current policy debates.

\subsubsection{Geographic clustering stability over twenty years}

A central finding of C.4 is that geographic clustering of district
characteristics is \emph{stable} across the three census cycles. Using
hierarchical clustering of state-level PP scores, C.4 identifies four
stable clusters: (1) compact inland states (PP $>$ 0.50); (2) moderate
coastal states (PP 0.40--0.50); (3) complex metropolitan states
(PP 0.35--0.40); and (4) irregular geographic states (PP $<$ 0.35).
Cluster membership changes for only 3 of 50 states across the
2000--2020 period, confirming that geographic characteristics dominate
algorithm performance in a stable, predictable manner.

\subsection{Implications for Decadal Planning}

The combined findings of C.3 and C.4 support a specific policy
conclusion: states that adopt recursive bisection in 2020 should expect
(a) moderately easier 2030 transitions, due to the $0.8$ pp structural
stability advantage; (b) stable algorithmic compactness performance
across the decade; and (c) a persistent but narrowing gap relative to
reform-era enacted plans. These properties make the algorithm suitable
for long-term planning contexts, not just single-census deployment.
